\scanpage{47}
\mspage{3}

L’ens des points fixes de $D'$ opérant sur
$\Sigma$ est \deleted{\unsure{indexé}} réduit à deux éléments,
un dans chacun des deux hémisphères
en lesquels $\Sigma_{\mathbb R}=P(V)$ découpe $\Sigma$ —
on a :
\[
(14)\qquad \pi_0(\Sigma\setminus\Sigma_{\mathbb R})\simeq
\begin{gathered}\underline\omega\\[-.5ex]\Vert\\[-.5ex]\Omega(\Sigma_{\mathbb R})\end{gathered}
\]
\margin{conv. de Stokes, provenant de l’orientation can. de $\Sigma$)}

donc les deux “pôles” seront indexés par
$\underline\omega$, et notés $\unsure{P|R}_\omega$
\[
(15)\qquad\left\{
\begin{aligned}
&P_\omega\qquad(\omega\in\underline\omega)\\
&P_\omega=\underbrace{\Sigma^{D^+}}_{\substack{\text{inv. sous}\\D^+\text{ op.}\\\text{sur }\Sigma}}
\cap\underbrace{\Sigma_\omega}_{\substack{\text{comp. conn.}\\\text{de }\Sigma\setminus\Sigma_{\mathbb R}\text{ indexée}\\\text{par }\omega}}
\end{aligned}\right.
\]

Si on exprime la définition de $\Sigma^{D^+}$ en
termes d’algèbre linéaire, on trouve que
ce sont les \uline{droites} \inserted{\uline{propres}} dans $V_{\mathbb C}$
\deleted{propres} \inserted{pour}
l’action de $D^+$. Mais considérons
la \unsure{double|seule} structure euclidienne sur $V_{\unsure{\mathbb C|\mathbb R}}$,
qui \deleted{\unsure{définit}} s’étend en une structure

\margin{46}
